\documentclass[11pt,spanish]{article}
\renewcommand{\sfdefault}{lmss}
\renewcommand{\familydefault}{\rmdefault}
\usepackage[utf8]{inputenc}
\usepackage{geometry}
\geometry{verbose,tmargin=2cm,bmargin=2cm,lmargin=2cm,rmargin=2cm,columnsep=1cm,marginparwidth=1cm}
\usepackage{fancyhdr}
\pagestyle{fancy}
\setlength{\headheight}{13.6pt}
\usepackage[skip=\bigskipamount]{parskip}
\usepackage{array}
\usepackage{microtype}
\usepackage{hyperref}
\hypersetup{hidelinks}
\usepackage{pifont}
\usepackage{titling}
\setlength{\droptitle}{-2cm}
\usepackage{xcolor}
\usepackage{enumitem}
\usepackage[most]{tcolorbox}
\usepackage{amsmath}
\usepackage{amssymb}
\usepackage{tikz}
\usetikzlibrary{automata,positioning,arrows.meta}

\definecolor{darkpastelblue}{rgb}{0.47, 0.62, 0.8}
\definecolor{amber}{rgb}{1.0, 0.75, 0.0}
\definecolor{rojo7a3c3d}{rgb}{0.478, 0.235, 0.239}
\definecolor{verde587246}{rgb}{0.345, 0.447, 0.275}
\definecolor{antiquefuchsia}{rgb}{0.57, 0.36, 0.51}
\definecolor{pastelgreen}{rgb}{0.82,0.92,0.78}
\definecolor{indiagreen}{rgb}{0.07,0.53,0.03}

\newcommand{\blank}[1][4cm]{\rule{#1}{0.4pt}}
\newcommand{\cuadro}{\(\square\)}
\newcommand{\cf}[2]{\colorbox{#1}{\textcolor{white}{\strut #2}}}
\ifrespuestas
\newcommand{\espaciosalida}{0.10cm}
\else
\newcommand{\espaciosalida}{0.60cm}
\fi

\newtcolorbox{ejercicio}[1]{
colback=antiquefuchsia!5!white,
colframe=antiquefuchsia!85!black,
fonttitle=\bfseries,
title={#1},
breakable}

\newtcolorbox{observacion}[1]{
colback=amber!5!white,
colframe=amber!85!black,
fonttitle=\bfseries,
title={#1},
breakable}

\newtcolorbox{respuesta}[1]{
colback=pastelgreen!35!white,
colframe=indiagreen!85!black,
fonttitle=\bfseries,
title={#1},
breakable}

\setlength{\parskip}{\bigskipamount}
\setlength{\parindent}{0pt}

\AtBeginDocument{
  \def\labelitemi{\(\star\)}
  \def\labelitemii{\(\ast\)}
  \def\labelitemiii{\(\bullet\)}
  \def\labelitemiv{\(\circ\)}
}

\usepackage{babel}
\addto\shorthandsspanish{\spanishdeactivate{~<>.}}
\deactivatequoting

\begin{document}

\lhead{PDS}
\rhead{Cuadernillo 03}
\title{Programación de Sistemas\\
Unidad 2: Análisis léxico\\ \bigskip
\cf{rojo7a3c3d}{\textbf{Cuadernillo de ejercicios 03}}\\
\bigskip
\Large\textbf{Expresiones regulares y autómatas}
\ifrespuestas\\[0.4em]\large Versión para docente\fi}
\author{Manuel Soto Romero \and Araceli Liliana Reyes Cabello}
\date{Semestre 2026-2 \\ Colegio de Ciencia y Tecnología UACM}

\maketitle
\renewcommand*{\thefootnote}{\arabic{footnote}}
\setcounter{footnote}{0}

Este cuadernillo acompaña la Nota 04. Su propósito es practicar la especificación de patrones mediante expresiones regulares, el reconocimiento con AFN y AFD, la construcción de subconjuntos y la minimización de autómatas. Las actividades avanzan de interpretar lenguajes pequeños a aplicar la ruta completa a los tokens del \textsc{IMP} básico.

\begin{observacion}{Observación 1. Límite de trabajo}
Responde únicamente con lo estudiado en la Nota 04. Utiliza las definiciones aprobadas \(\mathit{LOC}=\mathit{letra}^{+}\mathit{digito}^{*}\) y \(\mathit{NUM}=\mathit{digito}^{+}\). No necesitas implementar el analizador con \textsc{Lark}, coordinar todos los patrones ni diseñar recuperación de errores; esas tareas corresponden al tema siguiente.
\end{observacion}

\begin{observacion}{Observación 2. Ruta sugerida}
Los ejercicios \textbf{1, 2, 4, 5, 6, 8 y 10 son esenciales}: comprueban la interpretación de expresiones, el funcionamiento de los autómatas y las dos conversiones principales. Los ejercicios \textbf{3, 7 y 9 son de extensión}: exigen detectar especificaciones incorrectas o ejecutar procedimientos con un autómata distinto al ejemplo central. Si el tiempo es limitado, completa primero los esenciales.
\end{observacion}

\section*{\cf{verde587246}{1. Expresiones regulares y tokens de IMP}}

\begin{ejercicio}{Ejercicio 1. Cotejar cadenas con expresiones regulares \hfill {\small [Esencial]}}
Sea \(\Sigma=\{a,b\}\). Escribe \textbf{Sí} cuando la cadena pertenezca al lenguaje de la expresión y \textbf{No} cuando no pertenezca.

\begin{center}
\small
\renewcommand{\arraystretch}{1.3}
\begin{tabular}{|c|c|c|c|c|}
\hline
\textbf{Cadena} & \(a\mid b\) & \(ab\) & \(a^{*}\) & \(a(a\mid b)^{*}\) \tabularnewline
\hline
\(\varepsilon\) & & & & \tabularnewline[\espaciosalida]
\hline
\(a\) & & & & \tabularnewline[\espaciosalida]
\hline
\(b\) & & & & \tabularnewline[\espaciosalida]
\hline
\(ab\) & & & & \tabularnewline[\espaciosalida]
\hline
\(aa\) & & & & \tabularnewline[\espaciosalida]
\hline
\(ba\) & & & & \tabularnewline[\espaciosalida]
\hline
\end{tabular}
\end{center}

Explica por qué \(\varepsilon\) pertenece a \(L(a^{*})\), pero no a \(L(a(a\mid b)^{*})\).

\ifrespuestas\else
\blank[14cm]

\blank[14cm]
\fi
\end{ejercicio}

\ifrespuestas
\begin{respuesta}{Respuesta 1}
\begin{center}
\small
\renewcommand{\arraystretch}{1.15}
\begin{tabular}{|c|c|c|c|c|}
\hline
\textbf{Cadena} & \(a\mid b\) & \(ab\) & \(a^{*}\) & \(a(a\mid b)^{*}\) \tabularnewline
\hline
\(\varepsilon\) & No & No & Sí & No \tabularnewline
\hline
\(a\) & Sí & No & Sí & Sí \tabularnewline
\hline
\(b\) & Sí & No & No & No \tabularnewline
\hline
\(ab\) & No & Sí & No & Sí \tabularnewline
\hline
\(aa\) & No & No & Sí & Sí \tabularnewline
\hline
\(ba\) & No & No & No & No \tabularnewline
\hline
\end{tabular}
\end{center}

La cerradura \(a^{*}\) permite concatenar cero apariciones de \(a\), por lo que incluye \(\varepsilon\). En \(a(a\mid b)^{*}\), la primera \(a\) es obligatoria; la parte con cerradura puede ser vacía, pero la cadena completa no.
\end{respuesta}
\fi

\begin{ejercicio}{Ejercicio 2. Clasificar lexemas de IMP \hfill {\small [Esencial]}}
Clasifica cada cadena como \texttt{LOC}, \texttt{NUM} o \textbf{ninguna} cuando se considera como un solo lexema.

\begin{center}
\small
\renewcommand{\arraystretch}{1.25}
\begin{tabular}{|c|c||c|c|}
\hline
\textbf{Cadena} & \textbf{Clasificación} & \textbf{Cadena} & \textbf{Clasificación} \tabularnewline
\hline
\texttt{x} & & \texttt{20x} & \tabularnewline[\espaciosalida]
\hline
\texttt{total} & & \texttt{x2y} & \tabularnewline[\espaciosalida]
\hline
\texttt{x20} & & \texttt{x\_1} & \tabularnewline[\espaciosalida]
\hline
\texttt{0} & & \texttt{-5} & \tabularnewline[\espaciosalida]
\hline
\texttt{205} & & \(\varepsilon\) & \tabularnewline[\espaciosalida]
\hline
\end{tabular}
\end{center}

Justifica las clasificaciones de \texttt{x2y} y \texttt{-5} a partir de los patrones aprobados.

\ifrespuestas\else
\blank[14cm]

\blank[14cm]
\fi
\end{ejercicio}

\ifrespuestas
\begin{respuesta}{Respuesta 2}
\texttt{x}, \texttt{total} y \texttt{x20} son \texttt{LOC}; \texttt{0} y \texttt{205} son \texttt{NUM}. \texttt{20x}, \texttt{x2y}, \texttt{x\_1}, \texttt{-5} y \(\varepsilon\) no satisfacen ninguno de los dos patrones como un solo lexema.

En \texttt{x2y} aparece una letra después de comenzar la parte de dígitos, lo cual no permite \(\mathit{letra}^{+}\mathit{digito}^{*}\). En \texttt{-5}, el signo no es un dígito: durante el análisis léxico de un programa podría reconocerse como \texttt{MINUS} seguido de \texttt{NUM(5)}, pero la cadena completa no es un solo \texttt{NUM}.
\end{respuesta}
\fi

\begin{ejercicio}{Ejercicio 3. Encontrar un contraejemplo \hfill {\small [Extensión]}}
Cada propuesta intenta especificar un token de \textsc{IMP}, pero describe un lenguaje distinto al aprobado. Escribe una cadena que haga visible el error y corrige la expresión.

\begin{center}
\small
\renewcommand{\arraystretch}{1.35}
\begin{tabular}{|p{0.25\textwidth}|p{0.29\textwidth}|p{0.31\textwidth}|}
\hline
\textbf{Propuesta} & \textbf{Contraejemplo} & \textbf{Expresión corregida} \tabularnewline
\hline
\(\mathit{LOC}=\mathit{letra}\,\mathit{digito}^{*}\) & & \tabularnewline[0.9cm]
\hline
\(\mathit{LOC}=\mathit{letra}^{+}\mathit{digito}^{+}\) & & \tabularnewline[0.9cm]
\hline
\(\mathit{NUM}=\mathit{digito}^{*}\) & & \tabularnewline[0.9cm]
\hline
\end{tabular}
\end{center}
\end{ejercicio}

\ifrespuestas
\begin{respuesta}{Respuesta 3}
Ejemplos posibles:
\begin{itemize}[leftmargin=*]
    \item La primera propuesta rechaza \texttt{total} porque permite exactamente una letra antes de los dígitos. Debe ser \(\mathit{letra}^{+}\mathit{digito}^{*}\).
    \item La segunda propuesta rechaza \texttt{x} porque exige al menos un dígito. Debe ser \(\mathit{letra}^{+}\mathit{digito}^{*}\).
    \item La tercera propuesta acepta \(\varepsilon\), aunque \texttt{NUM} exige al menos un dígito. Debe ser \(\mathit{digito}^{+}\).
\end{itemize}
\end{respuesta}
\fi

\section*{\cf{verde587246}{2. Reconocimiento mediante AFN y AFD}}

\begin{ejercicio}{Ejercicio 4. Distinguir AFN y AFD \hfill {\small [Esencial]}}
Marca cada afirmación como \textbf{verdadera} o \textbf{falsa}. Corrige las afirmaciones falsas.

\begin{enumerate}[leftmargin=*]
    \item Un AFN puede incluir transiciones \(\varepsilon\). \hfill V \cuadro\quad F \cuadro
    \item En un AFD puede haber dos estados siguientes para el mismo estado y símbolo. \hfill V \cuadro\quad F \cuadro
    \item Un AFN acepta una cadena sólo cuando todos sus recorridos terminan en aceptación. \hfill V \cuadro\quad F \cuadro
    \item Una cadena determina un único recorrido en un AFD. \hfill V \cuadro\quad F \cuadro
    \item Para cada AFN existe un AFD que reconoce el mismo lenguaje. \hfill V \cuadro\quad F \cuadro
    \item Las transiciones \(\varepsilon\) consumen el siguiente símbolo de la entrada. \hfill V \cuadro\quad F \cuadro
\end{enumerate}

\ifrespuestas\else
\textbf{Correcciones:}

\blank[14cm]

\blank[14cm]

\blank[14cm]
\fi
\end{ejercicio}

\ifrespuestas
\begin{respuesta}{Respuesta 4}
1. Verdadera. 2. Falsa. 3. Falsa. 4. Verdadera. 5. Verdadera. 6. Falsa.

Correcciones:
\begin{itemize}[leftmargin=*]
    \item En un AFD existe exactamente un estado siguiente para cada par estado--símbolo.
    \item Un AFN acepta si existe al menos un recorrido que consume toda la entrada y termina en aceptación.
    \item Una transición \(\varepsilon\) cambia de estado sin consumir un símbolo de la entrada.
\end{itemize}
\end{respuesta}
\fi

\begin{ejercicio}{Ejercicio 5. Construir un AFD para \texttt{LOC} \hfill {\small [Esencial]}}
Construye un AFD para \(\mathit{LOC}=\mathit{letra}^{+}\mathit{digito}^{*}\). Utiliza el significado de los estados indicado y completa la tabla.

\begin{itemize}[leftmargin=*]
    \item \(q_0\): todavía no se consume ningún carácter.
    \item \(q_L\): ya se consumió una o más letras y todavía pueden aparecer letras o comenzar los dígitos.
    \item \(q_D\): ya comenzó la parte de dígitos; sólo pueden aparecer más dígitos.
    \item \(q_X\): estado muerto.
\end{itemize}

\begin{center}
\small
\renewcommand{\arraystretch}{1.3}
\begin{tabular}{|c|c|c|c|c|}
\hline
\textbf{Estado} & \textbf{letra} & \textbf{dígito} & \textbf{otro} & \textbf{¿Acepta?} \tabularnewline
\hline
\(q_0\) & & & & \tabularnewline[\espaciosalida]
\hline
\(q_L\) & & & & \tabularnewline[\espaciosalida]
\hline
\(q_D\) & & & & \tabularnewline[\espaciosalida]
\hline
\(q_X\) & & & & \tabularnewline[\espaciosalida]
\hline
\end{tabular}
\end{center}

Traza el recorrido de \texttt{x}, \texttt{x20}, \texttt{x2y} y \texttt{20}. Indica el estado final y si la cadena se acepta.

\ifrespuestas\else
\blank[14cm]

\blank[14cm]

\blank[14cm]
\fi
\end{ejercicio}

\ifrespuestas
\begin{respuesta}{Respuesta 5}
\begin{center}
\small
\renewcommand{\arraystretch}{1.15}
\begin{tabular}{|c|c|c|c|c|}
\hline
\textbf{Estado} & \textbf{letra} & \textbf{dígito} & \textbf{otro} & \textbf{¿Acepta?} \tabularnewline
\hline
\(q_0\) & \(q_L\) & \(q_X\) & \(q_X\) & No \tabularnewline
\hline
\(q_L\) & \(q_L\) & \(q_D\) & \(q_X\) & Sí \tabularnewline
\hline
\(q_D\) & \(q_X\) & \(q_D\) & \(q_X\) & Sí \tabularnewline
\hline
\(q_X\) & \(q_X\) & \(q_X\) & \(q_X\) & No \tabularnewline
\hline
\end{tabular}
\end{center}

Recorridos:
\begin{itemize}[leftmargin=*]
    \item \texttt{x}: \(q_0\to q_L\); acepta.
    \item \texttt{x20}: \(q_0\to q_L\to q_D\to q_D\); acepta.
    \item \texttt{x2y}: \(q_0\to q_L\to q_D\to q_X\); rechaza.
    \item \texttt{20}: \(q_0\to q_X\to q_X\); rechaza.
\end{itemize}
\end{respuesta}
\fi

\section*{\cf{verde587246}{3. Construcción de subconjuntos}}

\begin{ejercicio}{Ejercicio 6. Convertir el AFN de \(a\mid b\) \hfill {\small [Esencial]}}
Considera el AFN de Thompson utilizado en la Nota 04.

\begin{center}
\begin{tikzpicture}[shorten >=1pt,node distance=2.0cm,on grid,auto,>=Stealth]
\node[state,initial] (q0) {\(0\)};
\node[state,above right=of q0] (q1) {\(1\)};
\node[state,right=of q1] (q2) {\(2\)};
\node[state,below right=of q0] (q3) {\(3\)};
\node[state,right=of q3] (q4) {\(4\)};
\node[state,accepting,below right=of q2] (q5) {\(5\)};
\path[->]
  (q0) edge node[above left] {\(\varepsilon\)} (q1)
       edge node[below left] {\(\varepsilon\)} (q3)
  (q1) edge node {\(a\)} (q2)
  (q3) edge node {\(b\)} (q4)
  (q2) edge node[above right] {\(\varepsilon\)} (q5)
  (q4) edge node[below right] {\(\varepsilon\)} (q5);
\end{tikzpicture}
\end{center}

Completa los cálculos.

\begin{enumerate}[leftmargin=*]
    \item \(A=\operatorname{cerradura}_{\varepsilon}(\{0\})=\blank[5cm]\).
    \item \(\operatorname{mover}(A,a)=\blank[3cm]\) y su cerradura es \(B=\blank[4cm]\).
    \item \(\operatorname{mover}(A,b)=\blank[3cm]\) y su cerradura es \(C=\blank[4cm]\).
    \item Desde \(B\) o \(C\), leer \(a\) o \(b\) lleva a \(D=\blank[3cm]\).
\end{enumerate}

Completa la tabla del AFD. Marca como aceptación los conjuntos que contienen al estado \(5\).

\begin{center}
\small
\renewcommand{\arraystretch}{1.25}
\begin{tabular}{|c|c|c|c|}
\hline
\textbf{Estado} & \textbf{Con \(a\)} & \textbf{Con \(b\)} & \textbf{¿Acepta?} \tabularnewline
\hline
\(A\) & & & \tabularnewline[\espaciosalida]
\hline
\(B\) & & & \tabularnewline[\espaciosalida]
\hline
\(C\) & & & \tabularnewline[\espaciosalida]
\hline
\(D\) & & & \tabularnewline[\espaciosalida]
\hline
\end{tabular}
\end{center}

¿Cuántos subconjuntos contiene \(\mathcal{P}(Q)\) y cuántos de ellos son alcanzables en este AFD?

\ifrespuestas\else
\blank[14cm]
\fi
\end{ejercicio}

\ifrespuestas
\begin{respuesta}{Respuesta 6}
\begin{enumerate}[leftmargin=*]
    \item \(A=\{0,1,3\}\).
    \item \(\operatorname{mover}(A,a)=\{2\}\) y \(B=\{2,5\}\).
    \item \(\operatorname{mover}(A,b)=\{4\}\) y \(C=\{4,5\}\).
    \item \(D=\emptyset\).
\end{enumerate}

\begin{center}
\small
\begin{tabular}{|c|c|c|c|}
\hline
\textbf{Estado} & \textbf{Con \(a\)} & \textbf{Con \(b\)} & \textbf{¿Acepta?} \tabularnewline
\hline
\(A\) & \(B\) & \(C\) & No \tabularnewline
\hline
\(B\) & \(D\) & \(D\) & Sí \tabularnewline
\hline
\(C\) & \(D\) & \(D\) & Sí \tabularnewline
\hline
\(D\) & \(D\) & \(D\) & No \tabularnewline
\hline
\end{tabular}
\end{center}

Como \(|Q|=6\), el conjunto potencia contiene \(2^6=64\) subconjuntos. La construcción descubre sólo cuatro subconjuntos alcanzables: \(A\), \(B\), \(C\) y \(D\).
\end{respuesta}
\fi

\begin{ejercicio}{Ejercicio 7. Construir sólo los subconjuntos alcanzables \hfill {\small [Extensión]}}
El siguiente AFN sobre \(\Sigma=\{0,1\}\) procede del material de Teoría de la Computación. No tiene transiciones \(\varepsilon\).

\begin{center}
\begin{tikzpicture}[shorten >=1pt,node distance=3.0cm,on grid,auto,>=Stealth]
\node[state,initial] (q0) {\(q_0\)};
\node[state,right=of q0] (q1) {\(q_1\)};
\node[state,accepting,right=of q1] (q2) {\(q_2\)};
\path[->]
  (q0) edge[loop above] node {\(0,1\)} ()
       edge node {\(0\)} (q1)
  (q1) edge node {\(1\)} (q2);
\end{tikzpicture}
\end{center}

Parte del subconjunto inicial \(A=\{q_0\}\). Calcula las transiciones, agrega sólo los subconjuntos alcanzables y completa la tabla. Usa \(B=\{q_0,q_1\}\) y \(C=\{q_0,q_2\}\).

\begin{center}
\small
\renewcommand{\arraystretch}{1.3}
\begin{tabular}{|c|c|c|c|}
\hline
\textbf{Estado} & \textbf{Con 0} & \textbf{Con 1} & \textbf{¿Acepta?} \tabularnewline
\hline
\(A=\{q_0\}\) & & & \tabularnewline[\espaciosalida]
\hline
\(B=\{q_0,q_1\}\) & & & \tabularnewline[\espaciosalida]
\hline
\(C=\{q_0,q_2\}\) & & & \tabularnewline[\espaciosalida]
\hline
\end{tabular}
\end{center}

Traza \texttt{1101} en el AFD obtenido y describe el lenguaje reconocido.

\ifrespuestas\else
\blank[14cm]

\blank[14cm]
\fi
\end{ejercicio}

\ifrespuestas
\begin{respuesta}{Respuesta 7}
\begin{center}
\small
\begin{tabular}{|c|c|c|c|}
\hline
\textbf{Estado} & \textbf{Con 0} & \textbf{Con 1} & \textbf{¿Acepta?} \tabularnewline
\hline
\(A=\{q_0\}\) & \(B\) & \(A\) & No \tabularnewline
\hline
\(B=\{q_0,q_1\}\) & \(B\) & \(C\) & No \tabularnewline
\hline
\(C=\{q_0,q_2\}\) & \(B\) & \(A\) & Sí \tabularnewline
\hline
\end{tabular}
\end{center}

La traza es \(A\xrightarrow{1}A\xrightarrow{1}A\xrightarrow{0}B\xrightarrow{1}C\); por tanto, \texttt{1101} se acepta. El autómata reconoce las cadenas sobre \(\{0,1\}\) que terminan en \texttt{01}.
\end{respuesta}
\fi

\section*{\cf{verde587246}{4. Minimización de autómatas}}

\begin{ejercicio}{Ejercicio 8. Marcar pares distinguibles \hfill {\small [Esencial]}}
Retoma el AFD de cuatro estados del Ejercicio 6: \(A\) y \(D\) no son de aceptación; \(B\) y \(C\) sí. Escribe:

\begin{itemize}[leftmargin=*]
    \item \textbf{I} si el par se marca inicialmente porque sólo uno de sus estados es de aceptación;
    \item \textbf{T} si se marca después porque alguna transición lleva a un par ya marcado;
    \item \textbf{S} si permanece sin marca.
\end{itemize}

\begin{center}
\small
\renewcommand{\arraystretch}{1.25}
\begin{tabular}{|c|c|p{0.50\textwidth}|}
\hline
\textbf{Par} & \textbf{I/T/S} & \textbf{Símbolo o razón} \tabularnewline
\hline
\(\{A,B\}\) & & \tabularnewline[\espaciosalida]
\hline
\(\{A,C\}\) & & \tabularnewline[\espaciosalida]
\hline
\(\{A,D\}\) & & \tabularnewline[\espaciosalida]
\hline
\(\{B,C\}\) & & \tabularnewline[\espaciosalida]
\hline
\(\{B,D\}\) & & \tabularnewline[\espaciosalida]
\hline
\(\{C,D\}\) & & \tabularnewline[\espaciosalida]
\hline
\end{tabular}
\end{center}

Indica qué estados son equivalentes y completa el AFD mínimo.

\begin{center}
\small
\renewcommand{\arraystretch}{1.25}
\begin{tabular}{|c|c|c|c|}
\hline
\textbf{Estado mínimo} & \textbf{Estados originales} & \textbf{Con \(a\) o \(b\)} & \textbf{¿Acepta?} \tabularnewline
\hline
\(S\) & & & \tabularnewline[\espaciosalida]
\hline
\(F\) & & & \tabularnewline[\espaciosalida]
\hline
\(M\) & & & \tabularnewline[\espaciosalida]
\hline
\end{tabular}
\end{center}
\end{ejercicio}

\ifrespuestas
\begin{respuesta}{Respuesta 8}
\begin{center}
\small
\begin{tabular}{|c|c|p{0.50\textwidth}|}
\hline
\textbf{Par} & \textbf{I/T/S} & \textbf{Símbolo o razón} \tabularnewline
\hline
\(\{A,B\}\) & I & Sólo \(B\) acepta. \tabularnewline
\hline
\(\{A,C\}\) & I & Sólo \(C\) acepta. \tabularnewline
\hline
\(\{A,D\}\) & T & Con \(a\) llega a \(\{B,D\}\), que está marcado. \tabularnewline
\hline
\(\{B,C\}\) & S & Ambos llegan a \(D\) con \(a\) o \(b\). \tabularnewline
\hline
\(\{B,D\}\) & I & Sólo \(B\) acepta. \tabularnewline
\hline
\(\{C,D\}\) & I & Sólo \(C\) acepta. \tabularnewline
\hline
\end{tabular}
\end{center}

El único par equivalente es \(B\approx C\).

\begin{center}
\small
\begin{tabular}{|c|c|c|c|}
\hline
\textbf{Estado mínimo} & \textbf{Estados originales} & \textbf{Con \(a\) o \(b\)} & \textbf{¿Acepta?} \tabularnewline
\hline
\(S\) & \(\{A\}\) & \(F\) & No \tabularnewline
\hline
\(F\) & \(\{B,C\}\) & \(M\) & Sí \tabularnewline
\hline
\(M\) & \(\{D\}\) & \(M\) & No \tabularnewline
\hline
\end{tabular}
\end{center}
\end{respuesta}
\fi

\begin{ejercicio}{Ejercicio 9. Refinar una partición \hfill {\small [Extensión]}}
Minimiza el AFD de la tabla. El estado inicial es \(q_0\) y \(F=\{q_1,q_2,q_5\}\).

\begin{center}
\small
\renewcommand{\arraystretch}{1.2}
\begin{tabular}{|c|c|c|c|}
\hline
\textbf{Estado} & \textbf{Con \(a\)} & \textbf{Con \(b\)} & \textbf{¿Acepta?} \tabularnewline
\hline
\(q_0\) & \(q_1\) & \(q_2\) & No \tabularnewline
\hline
\(q_1\) & \(q_3\) & \(q_4\) & Sí \tabularnewline
\hline
\(q_2\) & \(q_4\) & \(q_3\) & Sí \tabularnewline
\hline
\(q_3\) & \(q_5\) & \(q_5\) & No \tabularnewline
\hline
\(q_4\) & \(q_5\) & \(q_5\) & No \tabularnewline
\hline
\(q_5\) & \(q_5\) & \(q_5\) & Sí \tabularnewline
\hline
\end{tabular}
\end{center}

Escribe la partición inicial y refínala hasta que no cambie.

\[
\Pi_0=\blank[10cm]
\]

\[
\Pi_f=\blank[10cm]
\]

Completa la tabla del AFD mínimo usando un estado nuevo por cada grupo de \(\Pi_f\).

\begin{center}
\small
\renewcommand{\arraystretch}{1.25}
\begin{tabular}{|c|c|c|c|c|}
\hline
\textbf{Nuevo} & \textbf{Estados originales} & \textbf{Con \(a\)} & \textbf{Con \(b\)} & \textbf{¿Acepta?} \tabularnewline
\hline
\(P_0\) & & & & \tabularnewline[\espaciosalida]
\hline
\(P_1\) & & & & \tabularnewline[\espaciosalida]
\hline
\(P_2\) & & & & \tabularnewline[\espaciosalida]
\hline
\(P_3\) & & & & \tabularnewline[\espaciosalida]
\hline
\end{tabular}
\end{center}
\end{ejercicio}

\ifrespuestas
\begin{respuesta}{Respuesta 9}
La partición inicial separa estados de aceptación y estados que no aceptan:
\[
\Pi_0=\big\{\{q_1,q_2,q_5\},\{q_0,q_3,q_4\}\big\}.
\]

Al refinar, \(q_5\) se separa de \(q_1,q_2\), y después \(q_0\) se separa de \(q_3,q_4\). La partición estable es
\[
\Pi_f=\big\{\{q_0\},\{q_1,q_2\},\{q_3,q_4\},\{q_5\}\big\}.
\]

\begin{center}
\small
\begin{tabular}{|c|c|c|c|c|}
\hline
\textbf{Nuevo} & \textbf{Estados originales} & \textbf{Con \(a\)} & \textbf{Con \(b\)} & \textbf{¿Acepta?} \tabularnewline
\hline
\(P_0\) & \(\{q_0\}\) & \(P_1\) & \(P_1\) & No \tabularnewline
\hline
\(P_1\) & \(\{q_1,q_2\}\) & \(P_2\) & \(P_2\) & Sí \tabularnewline
\hline
\(P_2\) & \(\{q_3,q_4\}\) & \(P_3\) & \(P_3\) & No \tabularnewline
\hline
\(P_3\) & \(\{q_5\}\) & \(P_3\) & \(P_3\) & Sí \tabularnewline
\hline
\end{tabular}
\end{center}
\end{respuesta}
\fi

\ifrespuestas\else\newpage\fi
\section*{\cf{verde587246}{5. Integración}}

\begin{ejercicio}{Ejercicio 10. Del patrón a los tokens de IMP \hfill {\small [Esencial]}}
Considera el siguiente fragmento de \textsc{IMP}:

\begin{center}
\texttt{if x20 <= 25 then x20 := x20 + 1 else skip}
\end{center}

Responde paso a paso.

\begin{enumerate}[leftmargin=*]
    \item Separa los lexemas con una barra vertical.

    \ifrespuestas\else
    \blank[14cm]
    \fi

    \item Escribe la secuencia conceptual de tokens. Conserva como atributos el texto de \texttt{LOC} y el valor de \texttt{NUM}.

    \ifrespuestas\else
    \blank[14cm]

    \blank[14cm]
    \fi

    \item Escribe el patrón que satisface \texttt{x20} y el que satisface \texttt{25}.

    \ifrespuestas\else
    \blank[14cm]
    \fi

    \item Ordena las representaciones para completar la ruta desde el patrón hasta el reconocedor reducido:

    \begin{center}
    \small
    \begin{tabular}{|c|c|c|c|}
    \hline
    AFD & expresión regular & AFD mínimo & AFN \tabularnewline
    \hline
    \end{tabular}
    \end{center}

    \[
    \text{patrón}\longrightarrow\blank[2.7cm]
    \longrightarrow\blank[2.7cm]
    \longrightarrow\blank[2.7cm]
    \longrightarrow\blank[2.7cm]
    \]

    \item Explica qué trabajo falta para pasar de reconocedores individuales al analizador léxico funcional del tema 2.3.

    \ifrespuestas\else
    \blank[14cm]

    \blank[14cm]
    \fi
\end{enumerate}
\end{ejercicio}

\ifrespuestas
\begin{respuesta}{Respuesta 10}
\begin{enumerate}[leftmargin=*]
    \item \texttt{if | x20 | <= | 25 | then | x20 | := | x20 | + | 1 | else | skip}.
    \item \texttt{IF LOC(x20) LE NUM(25) THEN LOC(x20) ASSIGN LOC(x20) PLUS NUM(1) ELSE SKIP}.
    \item \texttt{x20} satisface \(\mathit{letra}^{+}\mathit{digito}^{*}\); \texttt{25} satisface \(\mathit{digito}^{+}\).
    \item Patrón \(\longrightarrow\) expresión regular \(\longrightarrow\) AFN \(\longrightarrow\) AFD \(\longrightarrow\) AFD mínimo.
    \item Falta coordinar todas las reglas, resolver coincidencias como una palabra reservada frente a \texttt{LOC}, identificar el lexema adecuado dentro del flujo, omitir los espacios permitidos, reportar errores y expresar la especificación en \textsc{Lark}.
\end{enumerate}
\end{respuesta}
\fi

\end{document}
